% P10-4a, 14.09.2026: Mechanismus, Quellenreichweite und v4-Nachweis angeglichen.
% M07 / B-20, 11.09.2026: Quellenzeile aus meeting-4-ux-block-l.txt:33;
% F1-Ledger-Beleg hat abweichendes Präfix Leitung. Nicht stillschweigend
% korrigierter Laufbestand. Tabelleninhalte im historischen F1-Stand,
% keine späteren Erweiterungen aus dem heutigen Core übernehmen.

\begin{table}[htbp]
  \centering
  \small
  \begin{tabular}{>{\raggedright\arraybackslash}p{0.20\textwidth}>{\raggedright\arraybackslash}p{0.72\textwidth}}
    \hline
    \textbf{Station} & \textbf{Inhalt im dokumentierten Lauf} \\
    \hline
    Quelle & Aussage der Angehörigen-Vertreterin im Originaltranskript: „Angehörige sollen Besuche bei ihrer Bezugsperson vorab in der App ankündigen können, mit Datum und Uhrzeit, und die Pflegenden der Einrichtung sollen eine Übersicht der angekündigten Besuche sehen.`` \\
    Ledger & Ein quellgebundener Claim bündelt beide Teilaussagen („Angehörige sollen \ldots{} ankündigen können, und Pflegende \ldots{} sollen eine Übersicht der angekündigten Besuche sehen``) und enthält einen gespeicherten Belegtext; dessen Sprecherpräfix weicht vom Original ab (siehe Fließtext). \\
    Fachliche Ableitung & Die modellgestützte Ableitung zerlegt den gebündelten Claim in zwei normative Anforderungs-Aussagen mit vereinheitlichter Modalität („Angehörige \emph{müssen} \ldots{} ankündigen können`` / „Pflegende \emph{müssen} eine Übersicht \ldots{} sehen können``); der Claim-Bezug bleibt als Herkunftsangabe erhalten. \\
    Entscheidungsvorlage & Der Einordnungsplan stellt zwölf Operationen zur Entscheidung; für die Besuchsankündigung eine Neuaufnahme mit Begründung („Für Besuchsankündigungen existiert im Core keine passende Requirement-Entität``) samt Claim-Referenz. \\
    Übernahme & Der Mensch übernahm am Ingest-Gate acht Operationen und wies vier mit protokollierter Begründung ab; die deterministische Anwendung vergab die Core-Kennungen REQ-82 und REQ-83 und übernahm die Claim-Referenz samt Herkunftsangaben des Eingangs --- das zugehörige Feature und seine Arbeitspakete entstehen erst danach. \\
    Arbeitspakete und Themenzuordnung & Eine nachgelagerte modellgestützte Stufe schlug ein neues Feature („Besuchsankündigung und Besuchskoordination``) mit fachlicher Begründung vor; nach eigener menschlicher Freigabe (in diesem Lauf als dokumentierte Sammelfreigabe im UI) legte die Anwendung FC-15 und PBI-046/047 an, verband die PBIs über „deckt`` mit ihren Anforderungen und ordnete PBIs sowie Anforderungen dem Feature zu. \\
    Ergebnis & Die beiden Arbeitspakete wurden als GitHub-Issues \#44 („Besuche \ldots{} vorab ankündigen``) und \#45 („Übersicht der angekündigten Besuche \ldots``) real erzeugt; die Zuordnung PBI-zu-Issue wurde im Core gespeichert. \\
    \hline
  \end{tabular}
  \caption[Besuchsankündigung: von der Quelle zum Issue]{Ein
    fachlicher Inhalt über die Stationen des ausgeführten Falls F1
    (Abschnitt~\ref{sec:eval-fortschreibung}; Belegzuordnung in
    Anhang~\ref{anh:evaluation}). Die Quellenzeile folgt dem
    Originaltranskript, die übrigen Zeilen den Lauf-Artefakten.}
  \label{t:beispiel-besuchsplanung}
\end{table}
